% This is the chapter 'Modeling' to the "Hodgkin-Huxley neuron model V2.0 book"
% Last edited 2026.08.08

\ifx\magyar\undefined
\chapter[Abstract neuronal components]{Abstract neuronal components\label{sec:Modell-Components}}
\else
\chapter[Abstrakt neuron komponensek]{Abstrakt neuron komponensek\label{sec:Modell-Components}}
\fi

\iflatexml
\else
\minitoc
\fi



\ifx\magyar\undefined
\section[Ion channels]{Ion channels \label{sec:AbstractIon channels}}
\else
\section[Ioncsatornák]{Ioncsatornák \label{sec:AbstractIon channels}}
\fi

\ifx\magyar\undefined
Embedded in the wall of the membrane, a smaller amount of (controlled 
and uncontrolled) ion channels are scattered.
Between the membrane and the axon, there is an array of a large amount of uncontrolled ion channels (the Axon Initial Segment \gls{AIS}). 
The elastic wall of the membrane is covered with ion layers, and therefore, a potential difference is created there (the resting potential).
The membrane represents \textit{a confined space in which ions behave differently from in an infinite space}; furthermore, it significantly alters the properties of the few-nanometer-thick electrolyte layer near the membrane. In this layer, the concentration and potential change significantly with the distance from the membrane due to interactions between the membrane's ion layers and the electrolyte.
Ion packages (spikes) sent by other neurons can enter the internal volume through controlled inputs (synapse).

\else

A membrán falába ágyazva, nagyobb számú vezérelt és kisebb számú vezérlés nélküli ioncsatorna, a membrán és az axon között
pedig egy nagyon sok vezérletlen ioncsatornát tartalmazó ion tömb 
(az Axon Initial Segment \gls{AIS}) helyezkedik el.
A rugalmas membrán falát ionrétegek borítják, és ezért
potenciál különbség ("nyugalmi potenciál") keletkezik rajta.
\textit{A membrán olyan zárt teret képez, amelyben az ionok a szabad térben
ismerttől eltérő módon viselkednek}; továbbá, a membránon levő feszültség jelentősen megváltoztatja az elektrolitok tulajdonságait a 
membránhoz közel eső néhány nanométer vastagságú rétegben. 
Ebben a rétegben a koncentráció és a potenciál jelentősen függ
a membrántól való távolságtól; ami a membrán ionrétegei és az
elektrolit közötti kölcsönhatás következménye.
A más neuronok által küldött ion csomagok (spikes) a belső térfogatba
vezérelt bemeneteken (szinapszisokon)  át léphetnek be.

\fi

